\section*{अध्याय \ref{ch:b2:meiosis-heredity} --- अर्धसूत्री विभाजन, आनुवंशिक मिश्रण और वंशागति}

\begin{solution}{exo:b2:meiosis-heredity:1}
समसूत्री विभाजन: एक विभाजन, समजातों का कोई युग्मन नहीं, दो कोशिकाएँ, और
हर एक जनक से आनुवंशिक रूप से समरूप ($2n \to 2n$)। \omterm{def:b2:meiosis-heredity:meiosis}{अर्धसूत्री विभाजन}: एक
प्रतिकृतिकरण के बाद दो विभाजन, समजात युग्मित होकर जीन-विनिमय करते हैं,
चार कोशिकाएँ, \omterm{def:b2:meiosis-heredity:meiosis}{अगुणित} और सब आनुवंशिक रूप से भिन्न ($2n \to n$)।
\end{solution}

\begin{solution}{exo:b2:meiosis-heredity:2}
$2^{23} = 8.4\times 10^{6}$; $2^{4} = 16$; $2^{7} = 128$ --- और यह
जीन-विनिमय द्वारा इन सबको असीम गुणा करने से पहले।
\end{solution}

\begin{solution}{exo:b2:meiosis-heredity:3}
$Tt\times Tt$: \omterm{def:b2:meiosis-heredity:vocabulary}{जीन-प्ररूप} $1\,TT : 2\,Tt : 1\,tt$, \omterm{def:b2:meiosis-heredity:vocabulary}{लक्षण-प्ररूप} $3$ लंबे
: $1$ बौना। $Tt\times tt$: $1\,Tt : 1\,tt$, यानी $1 : 1$। उस लंबे पौधे
का परीक्षण संकरण किसी बौने से कराइए: सब लंबे आएँ तो $TT$, आधे बौने आएँ
तो $Tt$ (या उसे स्वपरागित कीजिए: $TT$ सत्य-प्रजनन करता है, $Tt$ $3 : 1$ पृथक होता है)।
\end{solution}

\begin{solution}{exo:b2:meiosis-heredity:4}
सहलग्न जीन एक ही गुणसूत्र पर होते हैं और युग्मकों में उनके जनकीय
\omterm{def:b2:meiosis-heredity:vocabulary}{युग्मविकल्पी} संयोजन अधिक होते हैं; \omterm{def:b2:meiosis-heredity:linkage}{पुनर्संयोजन
बारंबारता} पुनर्संयोजी युग्मकों का अंश है (किसी परीक्षण संकरण की
पुनर्संयोजी संतति); और एक \omterm{def:b2:meiosis-heredity:linkage}{सेंटीमॉर्गन} एक प्रतिशत
पुनर्संयोजन है। किसी लंबे गुणसूत्र पर दूर-दूर बैठे जीनों के बीच लगभग हर
\omterm{def:b2:meiosis-heredity:meiosis}{अर्धसूत्री विभाजन} में कम से कम एक जीन-विनिमय होता है, और अनेक
जीन-विनिमय संयोजनों को यादृच्छिक कर देते हैं, इसलिए $r$ $\tfrac12$ तक
पहुँच जाता है --- जो \omterm{thm:b2:meiosis-heredity:mixing}{स्वतंत्र अपव्यूहन} से अभेद्य है।
\end{solution}

\begin{solution}{exo:b2:meiosis-heredity:5}
कुल 556; प्रत्याशित $312.75 : 104.25 : 104.25 : 34.75$; $\chi^2 =
2.25^2/312.75 + 3.75^2/104.25 + 3.25^2/104.25 + 2.75^2/34.75 = 0.016
+ 0.135 + 0.101 + 0.218 = 0.47 < 7.81$: यानी $9 : 3 : 3 :
1$ के अनुरूप।
\end{solution}

\begin{solution}{exo:b2:meiosis-heredity:6}
$1 : 1 : 1 : 1$ के विरुद्ध (हर एक 250): $\chi^2 = (162^2 + 138^2 + 147^2
+ 153^2)/250 = 361 \gg 7.81$: सहलग्न। पुनर्संयोजी $103 + 97 = 200$
में से 1000: $r = 0.20$, \qty{20}{cM}। विकर्षण में $Ab/aB$ 400
$Ab$, 400 $aB$, 100 $AB$, 100 $ab$ देता है: वही दूरी, पर जनकीय
और पुनर्संयोजी वर्ग आपस में बदले हुए।
\end{solution}

\begin{solution}{exo:b2:meiosis-heredity:7}
$d = -\tfrac12\ln(1 - 2r)$: $r = 0.20$ देता है $-\tfrac12\ln 0.6 =
0.255$, यानी \qty{25.5}{cM}; $r = 0.45$ देता है $-\tfrac12\ln 0.1 = 1.15$,
यानी \qty{115}{cM}। $r = \tfrac12(1 - e^{-2d})$: \qty{30}{cM} देता है
$\tfrac12(1 - e^{-0.6}) = 0.226$; \qty{100}{cM} देता है $\tfrac12(1 -
e^{-2}) = 0.43$। \qty{10}{cM} से नीचे यह संशोधन उपेक्षणीय है;
\qty{30}{cM} से ऊपर वह अनिवार्य है, और \qty{50}{cM} से ऊपर $r$
संतृप्त हो जाता है, इसलिए कोई अकेला द्विबिंदु संकरण वह दूरी नाप ही नहीं
सकता --- दूर के जीन पास के जीनों की शृंखला जोड़कर मानचित्रित किए जाते हैं।
\end{solution}

\begin{solution}{exo:b2:meiosis-heredity:8}
(क) $X^wX^w \times X^+Y$: सारे पुत्र सफ़ेद ($X^wY$), सारी पुत्रियाँ लाल
($X^+X^w$)। (ख) $X^+X^w \times X^wY$: पुत्र आधे लाल, आधे सफ़ेद;
पुत्रियाँ आधी लाल ($X^+X^w$), आधी सफ़ेद ($X^wX^w$)। (ग) पुत्रियाँ
$X^+X^w$ $\times$ भाई $X^wY$: जैसा (ख) में --- हर लिंग के आधे
सफ़ेद, आधे लाल।
\end{solution}

\begin{solution}{exo:b2:meiosis-heredity:9}
शृंखला में दो एंजाइम वह रंजक बनाते हैं; बैंगनी के लिए दोनों स्थलों पर एक
\omterm{def:b2:meiosis-heredity:vocabulary}{प्रभावी} \omterm{def:b2:meiosis-heredity:vocabulary}{युग्मविकल्पी} चाहिए ($C\_P\_$)। वे वंशक्रम $CCpp$ और $ccPP$ हैं
(हर एक भिन्न कारण से सफ़ेद); संकर $CcPp$ बैंगनी है;
स्वपरागित होने पर वह $9\,C\_P\_$ बैंगनी बनाम $7$ सफ़ेद देता है ($3\,C\_pp + 3\,ccP\_ +
1\,ccpp$)। परीक्षण संकरण $CcPp\times ccpp$: $1$ बैंगनी ($CcPp$) : $3$
सफ़ेद।
\end{solution}

\begin{solution}{exo:b2:meiosis-heredity:10}
जनकीय $+++$, $abc$; दोहरे जीन-विनिमय $+b+$, $a+c$; जो जीन बदला वह $b$
है: क्रम $a$--$b$--$c$। $r_{ab} = (11 + 9 + 2)/1202
= 0.018$ (\qty{1.8}{cM}); $r_{bc} = (118 + 112 + 2)/1202 = 0.193$
(\qty{19.3}{cM})। प्रत्याशित दोहरे $0.018\times 0.193\times 1202 =
4.2$; प्रेक्षित 2: संपात $0.47$, व्यतिकरण $0.53$।
\end{solution}

\begin{solution}{exo:b2:meiosis-heredity:11}
द्वितीय-विभाजन ऐस्कस $300/1000 = \qty{30}{\%}$; दूरी $=
\tfrac12\times 30 = \qty{15}{cM}$। जीन और सेंट्रोमियर के बीच किसी
जीन-विनिमय में चार में से केवल दो क्रोमैटिड लगती हैं, इसलिए
द्वितीय-विभाजन पृथक्करण दिखाता कोई ऐस्कस दो पुनर्संयोजी और दो जनकीय
क्रोमैटिड ढोता है: \omterm{def:b2:meiosis-heredity:linkage}{पुनर्संयोजन बारंबारता} ऐसे ऐस्कसों की
बारंबारता की आधी है। $2 : 4 : 2$ के लिए: उस स्थल और सेंट्रोमियर के बीच
कोई जीन-विनिमय भीतरी दो क्रोमैटिड बदल देता है, इसलिए
\omterm{def:b2:meiosis-heredity:meiosis}{अर्धसूत्री विभाजन} I पर हर अर्ध-तर्कु एक $A$ और एक $a$
क्रोमैटिड ढोता है; फिर \omterm{def:b2:meiosis-heredity:meiosis}{अर्धसूत्री विभाजन} II उन्हें $A\,a \mid
a\,A$ (या उलटे) क्रम में अलग करता है, और अंतिम समसूत्री विभाजन हर बीजाणु
को दुगुना कर देता है: $2 : 4 : 2$।
\end{solution}

\begin{solution}{exo:b2:meiosis-heredity:12}
उसकी माता अनिवार्य वाहक है (उसे अपने पिता का एकमात्र X मिला था),
इसलिए उस स्त्री को वह \omterm{def:b2:meiosis-heredity:vocabulary}{युग्मविकल्पी} $\tfrac12$ प्रायिकता से मिला।
दो स्वस्थ पुत्र: वह वाहक हो तो दो स्वस्थ पुत्रों की प्रायिकता
$(\tfrac12)^2 = \tfrac14$ है, और अवाहक हो तो 1। बेज़ से:
$P(\text{वाहक}\mid\text{आँकड़े}) = (\tfrac12\times\tfrac14)/(\tfrac12
\times\tfrac14 + \tfrac12\times 1) = \tfrac{1/8}{5/8} = \tfrac15$।
\end{solution}

\begin{solution}{pb:b2:meiosis-heredity:1}
\textbf{1.} जनक $vg/vg\;+/+$ और $+/+\;e/e$; $F_1$ $+/vg\;+/e$;
युग्मक $+\,+$, $+\,e$, $vg\,+$, $vg\,e$, हर एक $\tfrac14$।
\textbf{2.} $900 : 300 : 300 : 100$।
\textbf{3.} $\chi^2 = 8^2/900 + 9^2/300 + 4^2/300 + 3^2/100 = 0.07 +
0.27 + 0.05 + 0.09 = 0.48 < 7.81$: \omterm{thm:b2:meiosis-heredity:mixing}{स्वतंत्र अपव्यूहन} चलता है।
\textbf{4.} दोनों स्थलों पर \omterm{def:b2:meiosis-heredity:vocabulary}{समयुग्मजी} वन्य सबका $\tfrac{1}{16}$ है,
और $\tfrac{9}{16}$ वन्य हैं: यानी $\tfrac19$।
\textbf{5.} वन्य, अवशेषी, आबनूसी, अवशेषी आबनूसी, हर एक $\tfrac14$।
\textbf{6.} \omterm{def:b2:meiosis-heredity:vocabulary}{लक्षण-प्ररूप} \omterm{def:b2:meiosis-heredity:vocabulary}{जीन-प्ररूप} और पर्यावरण का गुणनफल है
(ताप रंजक-एंजाइमों पर काम करता है): कोई \omterm{def:b2:meiosis-heredity:vocabulary}{जीन-प्ररूप} कोई
अनुक्रिया-मानक बताता है, कोई नियत लक्षण नहीं।
\textbf{7.} $b\,+/+\,vg$ (विकर्षण): $b$ एक जनक से $+_{vg}$ के साथ आया,
और $vg$ दूसरे से $+_b$ के साथ।
\textbf{8.} जनकीय: काली ($b\,+$) 965 और अवशेषी ($+\,vg$) 944;
पुनर्संयोजी: काली अवशेषी ($b\,vg$) 206 और वन्य ($+\,+$) 185।
वन्य-प्ररूप गुणसूत्र $+\,+$ $F_1$ में था ही नहीं; वह केवल किसी
जीन-विनिमय से ही बन सकता है।
\textbf{9.} $r = 391/2300 = 0.17$: \qty{17}{cM}।
\textbf{10.} $d = -\tfrac12\ln(1 - 0.34) = -\tfrac12\ln 0.66 = 0.21$:
\qty{21}{cM}।
\textbf{11.} नर \emph{Drosophila} में जीन-विनिमय होता ही नहीं:
नर के युग्मक केवल दोनों जनकीय गुणसूत्र ढोते हैं।
\textbf{12.} $F_1$ $b\,vg/+\,+$ (युग्मन): जनकीय वर्ग वन्य और
काली अवशेषी, हर एक \qty{41.5}{\%}; पुनर्संयोजी काली और
अवशेषी, हर एक \qty{8.5}{\%}।
\textbf{13.} $X^hX^+$: पुत्री को उसके पिता का एकमात्र X मिलता है, जो
$h$ ढोता है, और उसकी अप्रभावित माता से एक सामान्य X।
\textbf{14.} $\tfrac12$ (पुत्र) $\times$ $\tfrac12$ ($X^h$ पाना) $=
\tfrac14$।
\textbf{15.} $\tfrac12$: उसे माता के दोनों X में से कोई एक मिलता है;
पिता एक सामान्य X देता है।
\textbf{16.} हर संतान $\tfrac34$ प्रायिकता से अप्रभावित:
$(\tfrac34)^3 = \tfrac{27}{64} = 0.42$।
\textbf{17.} पुत्र को अपने पिता का Y मिलता है, उसका X कभी नहीं। कोई
स्त्री केवल $X^hX^h$ के रूप में प्रभावित होती है: इसके लिए प्रभावित
पिता और वाहक (या प्रभावित) माता चाहिए, जो विरल है।
\textbf{18.} पूर्व प्रायिकता $\tfrac12$। दो स्वस्थ पुत्रों की प्रायिकता
वह वाहक हो तो $\tfrac14$ है, न हो तो 1: पश्च प्रायिकता $(\tfrac12\times
\tfrac14)/(\tfrac12\times\tfrac14 + \tfrac12) = \tfrac15$।
\textbf{19.} पुत्र का एकमात्र X उसकी माता से आता है, और Y इनमें से कोई
जीन नहीं ढोता, इसलिए हर पुत्र एक मातृ युग्मक बिना किसी आवरण के दिखा
देता है: वह संकरण स्वयं ही अपना परीक्षण संकरण है।
\textbf{20.} जनकीय $+++$ (853), $ywm$ (833); दोहरे जीन-विनिमय
$+w+$ (2), $y+m$ (2); बदला $w$: क्रम $y$--$w$--$m$।
\textbf{21.} $r_{yw} = (24 + 26 + 2 + 2)/2000 = 0.027$; $r_{wm} = (128
+ 132 + 2 + 2)/2000 = 0.132$। मानचित्र: $y$ --- \qty{2.7}{cM} --- $w$ ---
\qty{13.2}{cM} --- $m$।
\textbf{22.} प्रत्याशित $0.027\times 0.132\times 2000 = 7.1$; प्रेक्षित
4: संपात $0.56$, व्यतिकरण $0.44$।
\textbf{23.} $y$--$m$ पुनर्संयोजी: $y++$, $+wm$, $yw+$, $++m$: $(24 +
26 + 128 + 132)/2000 = 0.155$, बनाम $0.027 + 0.132 = 0.159$: वे
चार दोहरे जीन-विनिमय जनकीय $y$--$m$ संयोजन लौटा देते हैं और छूट जाते
हैं, $2\times 4/2000 = 0.004$।
\textbf{24.} पुत्रियों को पिता का $y\,w\,m$ X और एक
मातृ X मिलता है: सब \omterm{def:b2:meiosis-heredity:vocabulary}{विषमयुग्मजी} हैं और \omterm{def:b2:meiosis-heredity:vocabulary}{लक्षण-प्ररूप} में वन्य, चाहे वे
कोई भी पुनर्संयोजी X ढोएँ; पुनर्संयोजन अदृश्य रहता है।
\textbf{25.} $y$ --- \qty{2.7}{cM} --- $w$ --- \qty{13.2}{cM} --- $m$;
संपात गुणांक $0.56$, व्यतिकरण $0.44$।
\end{solution}
